\begin{abstract}

Real-world sensor systems are subject to a broad range of hardware faults that degrade the quality
of collected data. However, Human Activity Recognition (HAR) research has largely focused on
clean data assumptions, and the research that considers non-clean data focuses on missing data
exclusively and provides imputation techniques to solve this problem. In this work, we introduce
a structured sensor corruption framework defined along three dimensions: corruption type, severity,
and granularity. We apply five physically motivated fault types: stochastic noise, dropout, bias,
gain, and drift, to the UCI HAR dataset and evaluate four models: KNN, Logistic Regression, SVM,
and LSTM. Our results show that bias and drift corruptions cause the most severe performance
degradation, while dropout produces less significant performance drops. The findings indicate the
importance of considering additional corruption types beyond dropout. We thus provide the
first baseline for how additional corruption types affect the HAR classification performance.
Our code is available in our public repository \cite{sensor-corruption-repo}.

\end{abstract}
